\documentclass[10pt,a4paper]{amsart}
\usepackage[margin=19mm]{geometry}
\usepackage{amsmath,amssymb,mathtools}
\usepackage[scr=rsfs,cal=euler]{mathalfa}
\usepackage{xfrac}
\usepackage[unicode,hidelinks,bookmarks=false]{hyperref}
\newcommand{\ie}{i.e.\ }
\newcommand{\cf}{cf.\ }
\newcommand{\resp}{respectively\ }
\newcommand{\Subsection}{\subsection}
\providecommand{\nobreakdash}{\nobreak\hbox{-}}
\setlength{\emergencystretch}{2em}
\multlinegap0pt
\usepackage{amssymb}
\newtheorem{proposition}{Proposition}
\newtheorem{lemma}{Lemma}
\newcommand{\MM}{{\mathfrak{M}}}
\newcommand{\NN}{{\mathfrak{N}}}
\title{Uniform Continuity in Distribution for Borel Transformations of Random Fields}
\author{Alexander I. Bufetov}
\date{Source: arXiv:2512.24441v1 (CC BY 4.0)}
\begin{document}
\maketitle

\noindent\emph{Source and licence.} Uniform Continuity in Distribution for Borel Transformations of Random Fields. Version arXiv:2512.24441v1 (CC BY 4.0), distributed by arXiv under the Creative Commons Attribution 4.0 International licence, http://creativecommons.org/licenses/by/4.0/, as recorded in the arXiv licence metadata for this exact version and shown on the abstract page https://arxiv.org/abs/2512.24441v1. Full licence notice: \texttt{licenses/arxiv-2512.24441-CC-BY-4.0.txt}.\par\smallskip

\noindent\textit{Editorial scope.} Counted unit: the article's proposition prop:3.1 (source0.tex lines 328--336). The article's complete original proof of it is delivered with it (source0.tex lines 338--354, one \texttt{proof} environment of 17 lines). No mathematics is written, repaired or re-derived: the statement and the proof are the article's own lines, in the article's own notation, and no retained line is substituted or re-flowed.  No further source-local context is needed: every cross-reference of every retained line resolves inside the two delivered spans.  Nothing was omitted from any retained span, so the package carries no omission record.  No retained line contains a citation, so the wrapper carries no bibliography and no bibliography entry.  No retained line contains a figure environment, an image-inclusion command or drawing code, so no image is delivered and the package declares no asset dependency.  Editorial formatting only, in addition: an AMS \texttt{amsart} preamble that restates the article's own declarations and macros used by the retained lines, namely the environments \texttt{proposition}, \texttt{lemma} on separate counters, together with the standard packages those lines need.  The retained environments print in this document as Proposition 1, Lemma 1 in order of appearance, whereas the article prints the same items with the numbering of its own full document.  The complete native source of the article as distributed in the arXiv e-print for this exact version is bundled unchanged as \texttt{sources/191895/source0.tex}.
	\begin{proposition}
		\label{prop:3.1}
		Let $M$ be a complete separable metric space, and let  $F \colon M \to C(\NN, \widetilde{\NN})$ be a Borel map.  If for any $l \in \mathbb{N}$ the 
		 finite-dimensional projection mapping
		\[
		m \mapsto \pi_l F(m)
		\]
		is uniformly continuous with respect to the L\'evy--Prokhorov metric, then the map $F$ is also uniformly continuous with respect to the L\'evy--Prokhorov metric.
	\end{proposition}

	\begin{proof}
		A  simple compactness argument  connects the distance between measures with
		the distances between their finite-dimensional distributions.
		
		\begin{lemma}
			\label{lem:3.2}
			Let $K \subset \MM_1(\mathbb{C}^{\mathbb{N}})$ be a compact set. For every $\varepsilon > 0$ there exist $l \in \mathbb{N}$
			and $\delta > 0$, depending only on $K$, such that
			if $\eta_1, \eta_2 \in K$ satisfy the inequality
			\[
			d_{\mathrm{LP}}(\operatorname{distr}^l_{\eta_1}([1, l]), \operatorname{distr}^l_{\eta_2}([1, l])) \le \delta,
			\]
			then  we have $$d_{\mathrm{LP}}(\eta_1, \eta_2) \le \varepsilon.$$
		\end{lemma}
		
		Lemma~\ref{lem:3.2} follows  from the Prokhorov Theorem and  directly implies Proposition \ref{prop:3.1}.
		\end{proof}

\end{document}
